\section{Estruturas de Dados}

\subsection{Fenwick Tree}

Suporta soma de prefixo e atualização pontual.

\textbf{Atualização:} $O(\log N)$

\textbf{Consulta de prefixo:} $O(\log N)$

\textbf{Soma em $[l,r]$:} $O(\log N)$

Usa índices de $1$ a $N$.

\begin{lstlisting}
int n;
vector<ll> bit;
void init (int size) {
    n = size;
    bit.assign(n + 1, 0);
}
void add (int i, ll x) {
    for (; i <= n; i += i & -i) bit[i] += x;
}
ll sum (int i) {
    ll res = 0;
    for (; i > 0; i -= i & -i) res += bit[i];
    return res;
}
ll query (int l, int r) {
    return sum(r) - sum(l - 1);
}
\end{lstlisting}


\subsection{Merge Sort Tree}

Permite consultas de ordem em intervalos usando
\texttt{lower\_bound} / \texttt{upper\_bound}.

\textbf{Construção:} $O(N\log N)$.

\textbf{Consulta:} $O(\log^2 N)$.

\begin{lstlisting}
int n;
vector<vector<int>> tree;
void build (const vector<int>& a, int p, int l, int r) {
    if (l == r) {
        tree[p] = {a[l]};
        return;
    }
    int m = (l + r) / 2;
    build(a, p * 2, l, m);
    build(a, p * 2 + 1, m + 1, r);
    merge(tree[p * 2].begin(), tree[p * 2].end(), 
          tree[p * 2 + 1].begin(), tree[p * 2 + 1].end(), back_inserter(tree[p]));
}
int query (int ql, int qr, int x, int p, int l, int r) {
    if (qr < l || r < ql) return 0;
    if (ql <= l && r <= qr)
        return lower_bound(tree[p].begin(), tree[p].end(), x) - tree[p].begin();
    int m = (l + r) / 2;
    return query(ql, qr, x, p * 2, l, m) + query(ql, qr, x, p * 2 + 1, m + 1, r);
}
\end{lstlisting}

O exemplo acima conta elementos $<x$ no intervalo $[l,r]$.

Para contar elementos $\leq x$, use \texttt{upper\_bound}.


\subsection{Li Chao Tree}

Mantém um conjunto de retas $y = mx + c$ e responde a consultas de menor valor em um ponto $x$ em $O(\log X)$, aceitando inserções arbitrárias.

\textbf{Inserção:} $O(\log X)$.

\textbf{Consulta:} $O(\log X)$.

\begin{lstlisting}
struct Line {
    ll m, c;
    ll eval(ll x) { return m * x + c; }
};
struct LiChao {
    int sz;
    vector<Line> tree;
    LiChao(int n) {
        sz = n;
        tree.assign(4 * sz, {0, INF}); // Use -INF se for para maximo
    }
    void add_line(int node, int l, int r, Line new_line) {
        int mid = (l + r) / 2;
        bool better_mid = new_line.eval(mid) < tree[node].eval(mid);
        bool better_l = new_line.eval(l) < tree[node].eval(l);
        if (better_mid) swap(tree[node], new_line);
        if (l == r) return;
        if (better_l != better_mid) 
            add_line(2 * node, l, mid, new_line);
        else
            add_line(2 * node + 1, mid + 1, r, new_line);
    }
    ll query(int node, int l, int r, ll x) {
        int mid = (l + r) / 2;
        ll res = tree[node].eval(x);
        if (l == r) return res;
        if (x <= mid) return min(res, query(2 * node, l, mid, x));
        else return min(res, query(2 * node + 1, mid + 1, r, x));
    }
};
\end{lstlisting}


\subsection{Persistent Segment Tree}

Mantém múltiplas versões de uma árvore de segmentos. Cada atualização cria apenas um novo caminho de nós ($O(\log N)$ novos nós por operação).

\textbf{Complexidade:} $O(\log N)$ em tempo e espaço por versão criada ou consulta.

\begin{lstlisting}
const int MAXN = 2e5 + 5;
const int MAXNODES = MAXN * 40; // Ajuste conforme os limites de memoria

struct Node {
    int val;
    int left, right;
} pool[MAXNODES];

int root_ver[MAXN];
int node_cnt = 0;

int update(int prev_node, int l, int r, int idx, int val) {
    int cur_node = ++node_cnt;
    pool[cur_node] = pool[prev_node];
    pool[cur_node].val += val;
    if (l == r) return cur_node;
    int mid = (l + r) / 2;
    if (idx <= mid)
        pool[cur_node].left = update(pool[prev_node].left, l, mid, idx, val);
    else 
        pool[cur_node].right = update(pool[prev_node].right, mid + 1, r, idx, val);
    return cur_node;
}

int query(int node, int l, int r, int ql, int qr) {
    if (!node || ql > r || qr < l) return 0;
    if (ql <= l && r <= qr) return pool[node].val;
    int mid = (l + r) / 2;
    return query(pool[node].left, l, mid, ql, qr) + 
           query(pool[node].right, mid + 1, r, ql, qr);
}
\end{lstlisting}


\subsection{Treap Implícita}

Mantém uma sequência dinâmica de elementos permitindo inserções, remoções e inversões de intervalos (\textit{reverse}) em $O(\log N)$.

\textbf{Complexidade:} $O(\log N)$ por operação de \texttt{split}, \texttt{merge} e modificação de intervalo.

\begin{lstlisting}
mt19937 rng(1337);

struct Node {
    int val, prior, sz;
    bool rev;
    Node *l, *r;
    Node(int v) : val(v), prior(rng()), sz(1), rev(false), l(nullptr), r(nullptr) {}
};

using pnode = Node*;

int get_sz(pnode t) { return t ? t->sz : 0; }
void update(pnode t) {
    if (t) t->sz = 1 + get_sz(t->l) + get_sz(t->r);
}
void push(pnode t) {
    if (t && t->rev) {
        t->rev = false;
        swap(t->l, t->r);
        if (t->l) t->l->rev ^= 1;
        if (t->r) t->r->rev ^= 1;
    }
}
// Divide a Treap t em L (primeiros 'key' elementos) e R (o restante)
void split(pnode t, pnode &l, pnode &r, int key, int add = 0) {
    if (!t) return void(l = r = nullptr);
    push(t);
    int cur_key = add + get_sz(t->l);
    if (key <= cur_key) {
        split(t->l, l, t->l, key, add);
        r = t;
    } else {
        split(t->r, t->r, r, key, cur_key + 1);
        l = t;
    }
    update(t);
}
// Une as Treaps l e r (assume que todos em l vem antes de r)
void merge(pnode &t, pnode l, pnode r) {
    push(l); push(r);
    if (!l || !r) t = l ? l : r;
    else if (l->prior > r->prior) {
        merge(l->r, l->r, r);
        t = l;
    } else {
        merge(r->l, l, r->l);
        t = r;
    }
    update(t);
}
// Inverte o subvetor no intervalo [l, r] (0-indexed)
void reverse_range(pnode &t, int l, int r) {
    pnode t1, t2, t3;
    split(t, t1, t2, l);
    split(t2, t2, t3, r - l + 1);
    if (t2) t2->rev ^= 1;
    merge(t2, t2, t3);
    merge(t, t1, t2);
}
\end{lstlisting}
